\documentclass[11pt]{article}

\usepackage[T1]{fontenc}
\usepackage{lmodern}
\usepackage[margin=1in]{geometry}
\usepackage{amsmath,amssymb,amsthm,mathtools}
\usepackage{booktabs}
\usepackage{enumitem}
\usepackage{float}
\usepackage{microtype}
\usepackage[hidelinks]{hyperref}

\newtheorem{theorem}{Theorem}
\newtheorem{lemma}[theorem]{Lemma}
\newtheorem{claim}[theorem]{Claim}
\newtheorem{proposition}[theorem]{Proposition}
\newtheorem{corollary}[theorem]{Corollary}
\theoremstyle{definition}
\newtheorem{definition}[theorem]{Definition}
\theoremstyle{remark}
\newtheorem{remark}[theorem]{Remark}

\newcommand{\A}{\mathcal A}
\newcommand{\C}{\mathcal C}
\newcommand{\F}{\mathcal F}
\newcommand{\Hh}{\mathcal H}
\newcommand{\I}{\mathcal I}
\newcommand{\N}{\mathcal N}
\newcommand{\Pp}{\mathcal P}
\newcommand{\Sset}{\mathcal S}
\newcommand{\OPT}{\operatorname{OPT}}
\newcommand{\Obj}{\operatorname{Obj}}
\newcommand{\err}{\operatorname{err}}
\newcommand{\dist}{\operatorname{dist}}
\newcommand{\Agree}{\mathsf{Agree}}
\newcommand{\maj}{\operatorname{maj}}
\newcommand{\rec}{\operatorname{rec}}
\newcommand{\eps}{\varepsilon}
\newcommand{\1}{\mathbf 1}

\allowdisplaybreaks

\hypersetup{
  pdftitle={A Strictly Better-than-Two Competitive Algorithm for Online Consensus Clustering},
  pdfsubject={Online consensus clustering with polylogarithmic per-vertex recourse}
}

\title{A Strictly Better-than-Two Competitive Algorithm\\
for Online Consensus Clustering}
\author{}
\date{}

\begin{document}
\maketitle

\section{A better-than-two competitive algorithm for Online Consensus Clustering}

In this section, we give a randomized online algorithm for consensus
clustering whose competitive ratio is strictly smaller than $2$ with high
probability.  The algorithm has polylogarithmic recourse per vertex.  The
result combines two ingredients.  The first ingredient is the sampled
cluster-fitting framework of Chakraborty, Chatterjee, Das, and
Nguyen~\cite{ChakrabortyEtAl2026}.  The second ingredient is a rebuilding
strategy based on monotone history scores.  The history score of a candidate
is the maximum consensus cost that the candidate has incurred so far.  This
maximum, rather than the current cost, is used to decide when the maintained
solution must be rebuilt.

The cluster-fitting framework is crucial.  Sampling only input clusterings
cannot give a competitive ratio bounded away from $2$ when the number $m$ of
input clusterings is arbitrary.  In contrast, cluster fitting also generates
candidates from triples of sampled input clusterings.  A suitable triple
candidate has cost at most $3\OPT/2$.  We prove that this candidate remains
good, relative to the optimum at the end of the current prefix, throughout
its entire history.  This is precisely the stability property needed by the
rebuilding argument.

We use the deterministic Online Agreement algorithm of Cohen-Addad,
Lattanzi, Maggiori, and Parotsidis~\cite{CohenAddadEtAl2022} as a black box.
Their result gives absolute constants $\eta\ge1$ and $q_{\mathrm A}>0$ with
the following property.

\begin{theorem}[Online Agreement; black-box form]\label{thm:online-agreement}
There is a deterministic algorithm $\Agree$ for unweighted complete signed
correlation clustering under vertex arrivals such that, at every prefix,
\[
  \err_H(\Agree(H))\le \eta\min_{\C}\err_H(\C).
\]
where $\err_H$ represents the number of correlation clustering disagreements made on the signed graph $H$.

Moreover, every vertex changes its cluster identifier at most
\[
  q_{\mathrm A}\lceil\log_2 n\rceil
\]
times during the entire execution.  Here $\eta$ and $q_{\mathrm A}$ are
absolute constants.
\end{theorem}

The published analysis states $\eta$ as an unspecified absolute constant.
Consequently, the improvement below $2$ obtained in this section is an
absolute constant, although its numerical value may be small.

We now state the main result.  An arrival sequence is called
\emph{oblivious} if the input clusterings and the vertex-arrival order are
fixed independently of the random samples used by the algorithm.

\begin{theorem}[Main theorem]\label{thm:main}
Fix a failure probability $\zeta\in(0,1)$.  Against every oblivious instance
of online consensus clustering, there is a randomized online algorithm such
that, with probability at least $1-\zeta$, simultaneously for every prefix
$t$,
\begin{equation}\label{eq:main-ratio}
  \Obj_t(\C_t)
  \le
  \left(2-\frac{1}{288(\eta+1)}\right)\OPT_t.
\end{equation}
Here $\eta$ is the deterministic competitive ratio in
Theorem~\ref{thm:online-agreement}.  On every outcome of the random samples,
every vertex has
\[
  O_{\eta,q_{\mathrm A}}\!\left(
     \log^4\!\frac{mn}{\zeta}
  \right)
\]
recourse.  In particular, the competitive ratio is strictly smaller than
$2$, and the recourse bound is worst-case rather than merely a
high-probability bound.
\end{theorem}

\subsection{Problem definition and notation}

In this subsection, we define the online consensus clustering problem and
all the notation used in the proof.

Let
\[
  V=\{v_1,v_2,\ldots,v_n\}
\]
be the ground set of vertices.  The vertices arrive in the fixed order
$v_1,v_2,\ldots,v_n$, and
\[
  V_t:=\{v_1,v_2,\ldots,v_t\}
\]
denotes the set of vertices that have arrived by time $t$.  We are given
$m$ input clusterings
\[
  \Pp_1,\Pp_2,\ldots,\Pp_m
\]
of $V$.  When $v_t$ arrives, the algorithm learns the cluster containing
$v_t$ in each input clustering.  Equivalently, for every $u\in V_{t-1}$ and
every $i\in[m]$, the algorithm learns whether $u$ and $v_t$ are together or
separated in $\Pp_i$.

For a clustering $\C$ and two vertices $u,v$, write
$u\sim_{\C}v$ when $u$ and $v$ belong to the same cluster of $\C$, and write
$u\not\sim_{\C}v$ otherwise.  For two clusterings $\C$ and $\C'$ of the same
vertex set $W$, define their pair-disagreement distance by
\begin{equation}\label{eq:partition-distance}
 \dist_W(\C,\C')
 :=
 \left|
 \left\{
   \{u,v\}\subseteq W:
   \1[u\sim_{\C}v]\ne\1[u\sim_{\C'}v]
 \right\}
 \right|.
\end{equation}
Thus, $\dist_W(\C,\C')$ is the number of unordered vertex pairs on
which $\C$ and $\C'$ disagree.  This is the Hamming distance between the
two pair-relation vectors, and therefore it satisfies the triangle
inequality.

Let
\[
  \Pp_i^t:=\Pp_i|_{V_t}
\]
be the restriction of the $i$-th input clustering to the arrived vertices.
For a clustering $\C$ of $V_t$, its unnormalized consensus objective is
\begin{equation}\label{eq:consensus-objective}
  \Obj_t(\C)
  :=\sum_{i=1}^m\dist_{V_t}(\C,\Pp_i^t).
\end{equation}
We use the unnormalized objective only to keep all costs integral.  Dividing
by $m$ gives the normalized metric-correlation-clustering objective and does
not change any competitive ratio.  Define
\begin{equation}\label{eq:opt-definition}
  \OPT_t:=\min_{\C}\Obj_t(\C),
\end{equation}
where the minimum is over all clusterings of $V_t$.

The sequence $\OPT_1,\OPT_2,\ldots,\OPT_n$ is nondecreasing.  Indeed, the
restriction to $V_t$ of any clustering of $V_{t+1}$ has no larger cost than
the original clustering because all pairs containing $v_{t+1}$ are removed.
Applying this observation to an optimum at time $t+1$ gives
\begin{equation}\label{eq:opt-monotone}
  \OPT_t\le\OPT_{t+1}.
\end{equation}

\paragraph{Signed correlation clustering.}
Let $H$ be a complete graph in which every edge is labeled either positive
or negative.  For a clustering $\C$, define
\begin{equation}\label{eq:signed-error}
 \err_H(\C)
 :=
 \bigl|
   \{\{u,v\}:\{u,v\}\text{ is positive and }u\not\sim_{\C}v\}
 \bigr|
 +
 \bigl|
   \{\{u,v\}:\{u,v\}\text{ is negative and }u\sim_{\C}v\}
 \bigr|.
\end{equation}

\paragraph{Recourse.}
Every output cluster carries a persistent identifier.  An old vertex incurs
one unit of recourse whenever its cluster identifier changes between two
consecutive time steps.  The first identifier assigned to a newly arriving
vertex is not counted.  This is the standard per-vertex recourse convention.

If $m=1$, outputting $\Pp_1^t$ is optimal and has zero recourse.  If $m=2$,
the triangle inequality gives, for every clustering $\C$,
\[
 \dist(\Pp_1^t,\Pp_2^t)
 \le
 \dist(\Pp_1^t,\C)+\dist(\C,\Pp_2^t),
\]
and $\Pp_1^t$ attains equality.  Thus, outputting $\Pp_1^t$ is again optimal
and has zero recourse.  In the rest of the proof, we assume $m\ge3$.

\subsection{The candidate clusterings}

In this subsection, we describe the candidate clusterings maintained by the
algorithm.  There are two types of candidates.

\paragraph{Input candidates.}
For every sampled index $i\in[m]$, the corresponding input candidate at time
$t$ is simply
\[
  \A_i(t):=\Pp_i^t.
\]
This candidate has zero internal recourse because it is the restriction of a
fixed partition of the final vertex set.

\paragraph{Triple candidates.}
Fix a triple $\tau=(i,j,k)$ of sampled input indices.  At time $t$, form the
complete signed graph $H_\tau(t)$ on $V_t$ as follows.  For every pair
$\{u,v\}\subseteq V_t$, label the edge positive if at least two of
$\Pp_i^t,\Pp_j^t,\Pp_k^t$ put $u$ and $v$ together, and label it negative if
at least two of them separate $u$ and $v$.  Thus, the edge label is the
majority pair relation of the three input clusterings.

The graphs
\[
  H_\tau(1),H_\tau(2),\ldots,H_\tau(n)
\]
form a valid vertex-arrival instance of signed correlation clustering.  We
run one deterministic copy of $\Agree$ on this sequence and denote its
clustering at time $t$ by
\[
  \A_\tau(t):=\Agree(H_\tau(t)).
\]
Theorem~\ref{thm:online-agreement} gives
\begin{equation}\label{eq:triple-cc-guarantee}
  \err_{H_\tau(t)}(\A_\tau(t))
  \le
  \eta\min_{\C}\err_{H_\tau(t)}(\C)
\end{equation}
at every prefix, deterministically.

For every candidate $a$, where $a$ is either an input index or a triple,
define its current full consensus cost by
\begin{equation}\label{eq:current-candidate-score}
  f_a(t):=\Obj_t(\A_a(t)).
\end{equation}
The current cost of a triple candidate need not be nondecreasing because
$\Agree$ may move old vertices and thereby decrease its cost.  We therefore
define its \emph{history score}
\begin{equation}\label{eq:history-score}
  G_a(t):=\max_{1\le u\le t}f_a(u).
\end{equation}
Every $G_a(t)$ is nondecreasing in $t$.  For an input candidate $i$, the
quantity $f_i(t)$ is itself nondecreasing, because each distance
$\dist_{V_t}(\Pp_i^t,\Pp_h^t)$ only gains newly revealed pairs.  Hence
\begin{equation}\label{eq:input-history-current}
  G_i(t)=f_i(t).
\end{equation}

\subsection{A sampled cluster-fitting lemma}

In this subsection, we give a self-contained specialization of the sampled
cluster-fitting argument from~\cite{ChakrabortyEtAl2026}.  The lemma is
stated for one fixed prefix.  Later, we amplify its probability and apply it
simultaneously to all prefixes.

Fix a time $t$ and suppress the superscript $t$ throughout this subsection.
Let $\F^*$ be an optimal consensus clustering, so that
\[
  \sum_{i=1}^m\dist(\Pp_i,\F^*)=\OPT.
\]
Define
\begin{equation}\label{eq:avg-definition}
  \mathsf{AVG}:=\frac{\OPT}{m}.
\end{equation}
For each input clustering $\Pp_i$, let
\begin{equation}\label{eq:error-set}
  I_i
  :=
  \left\{
    \{u,v\}:\Pp_i\text{ and }\F^*
    \text{ disagree on the pair }\{u,v\}
  \right\}.
\end{equation}
Then
\begin{equation}\label{eq:error-size}
  |I_i|=\dist(\Pp_i,\F^*)
  \qquad\text{and}\qquad
  \sum_{i=1}^m|I_i|=\OPT=m\mathsf{AVG}.
\end{equation}

For any two input clusterings, their disagreement set is the symmetric
difference $I_i\triangle I_j$.  Consequently,
\begin{equation}\label{eq:input-distance-error-sets}
 \dist(\Pp_i,\Pp_j)
 =|I_i|+|I_j|-2|I_i\cap I_j|.
\end{equation}
Summing~\eqref{eq:input-distance-error-sets} over $j$ gives the following
identity for the cost of an input candidate:
\begin{equation}\label{eq:input-candidate-cost-identity}
 \Obj(\Pp_i)
 =\OPT+m|I_i|-2\sum_{j=1}^m|I_i\cap I_j|.
\end{equation}

\subsubsection{The sparse-triple candidate}

We first analyze a triple whose three error sets have small pairwise
intersections.

\begin{lemma}[Sparse-triple fitting]\label{lem:sparse-triple}
Suppose a triple $\tau=(i,j,k)$ satisfies
\begin{equation}\label{eq:sparse-triple-condition}
 |I_i\cap I_j|\le\alpha\mathsf{AVG},\qquad
 |I_j\cap I_k|\le\alpha\mathsf{AVG},\qquad
 |I_k\cap I_i|\le\alpha\mathsf{AVG}.
\end{equation}
Let $\A_\tau$ be any $\eta$-approximate correlation clustering of the
majority signed graph $H_\tau$.  Then
\begin{equation}\label{eq:sparse-triple-cost}
 \Obj(\A_\tau)
 \le
 \bigl(1+3(\eta+1)\alpha\bigr)\OPT.
\end{equation}
\end{lemma}

\begin{proof}
Let $B$ be the set of pairs on which the majority relation of
$\Pp_i,\Pp_j,\Pp_k$ disagrees with $\F^*$.  A binary majority differs from
$\F^*$ exactly when at least two of the three input relations differ from
$\F^*$.  Therefore,
\begin{equation}\label{eq:bad-set-union}
 B=(I_i\cap I_j)\cup(I_j\cap I_k)\cup(I_k\cap I_i).
\end{equation}
By~\eqref{eq:sparse-triple-condition},
\begin{equation}\label{eq:bad-set-bound}
 |B|\le3\alpha\mathsf{AVG}.
\end{equation}

The clustering $\F^*$ makes exactly $|B|$ mistakes on the signed graph
$H_\tau$.  Hence the optimum correlation-clustering error on $H_\tau$ is at
most $|B|$, and
\[
  \err_{H_\tau}(\A_\tau)\le\eta|B|.
\]
The pair-relation vectors of $\A_\tau$ and $\F^*$ can be compared through
the majority sign vector.  The Hamming triangle inequality gives
\begin{equation}\label{eq:distance-through-majority}
 \dist(\A_\tau,\F^*)
 \le
 \err_{H_\tau}(\A_\tau)+\err_{H_\tau}(\F^*)
 \le(\eta+1)|B|.
\end{equation}
Using the triangle inequality once more, now between each input clustering
and $\A_\tau$, we obtain
\begin{align*}
 \Obj(\A_\tau)
 &=\sum_{h=1}^m\dist(\Pp_h,\A_\tau)\\
 &\le\sum_{h=1}^m
       \bigl(\dist(\Pp_h,\F^*)+\dist(\F^*,\A_\tau)\bigr)\\
 &\le\OPT+m(\eta+1)|B|\\
 &\le\bigl(1+3(\eta+1)\alpha\bigr)\OPT,
\end{align*}
where the final inequality uses
$\mathsf{AVG}=\OPT/m$ and~\eqref{eq:bad-set-bound}.
\end{proof}

\subsubsection{Bounds for ordered input candidates}

Relabel the input clusterings for the analysis so that
\begin{equation}\label{eq:ordered-errors}
  |I_1|\le|I_2|\le\cdots\le|I_m|.
\end{equation}
The algorithm does not know or use this ordering.

Fix parameters
\[
  \alpha>0,\qquad 0<\beta<1,\qquad c>1,\qquad s>1,
  \qquad g>s,
\]
such that $cs-s-2c>0$.  Let $t_0$ be the largest index such that
\begin{equation}\label{eq:low-inputs}
  |I_{t_0}|\le(1-\beta)\mathsf{AVG}.
\end{equation}
If there is no such index, set $t_0=0$.

\begin{lemma}[Ordered-error bounds]\label{lem:ordered-error-bounds}
Suppose $t_0<m/g$.
\begin{enumerate}[label=(\roman*),leftmargin=2.5em]
\item For every $h$ satisfying $t_0<h\le m/s$,
\begin{equation}\label{eq:h-error-bound}
 |I_h|
 \le
 \left(
  1+\frac{\beta(g-s)+s}{g(s-1)}
 \right)\mathsf{AVG}.
\end{equation}
\item For every $k$ satisfying
\[
 t_0<k\le\frac{2c+s}{cs}m,
\]
we have
\begin{equation}\label{eq:k-error-bound}
 |I_k|
 \le
 \left(
 1+
 \frac{\beta\bigl(g(2c+s)-sc\bigr)+sc}
      {g(cs-s-2c)}
 \right)\mathsf{AVG}.
\end{equation}
\end{enumerate}
\end{lemma}

\begin{proof}
Fix an index $a>t_0$.  By the ordering and the definition of $t_0$,
the indices $t_0+1,\ldots,a-1$ have error larger than
$(1-\beta)\mathsf{AVG}$, while the indices $a,a+1,\ldots,m$ have error at
least $|I_a|$.  Therefore,
\begin{align*}
 m\mathsf{AVG}
 =\sum_{i=1}^m|I_i|
 &\ge
 (a-t_0-1)(1-\beta)\mathsf{AVG}
 +(m-a+1)|I_a|.
\end{align*}
Rearranging gives
\begin{equation}\label{eq:generic-order-bound}
 \frac{|I_a|}{\mathsf{AVG}}
 \le
 1+
 \frac{\beta(a-1)+(1-\beta)t_0}{m-a+1}.
\end{equation}

If $a=h\le m/s$, then
\[
 \frac{h-1}{m-h+1}\le\frac1{s-1}
 \qquad\text{and}\qquad
 \frac{t_0}{m-h+1}\le\frac{s}{g(s-1)}.
\]
Substituting these two inequalities into
\eqref{eq:generic-order-bound} proves~\eqref{eq:h-error-bound}.

For the second part, put
\[
 \lambda:=\frac{2c+s}{cs}.
\]
The assumption $cs-s-2c>0$ is exactly the assertion $\lambda<1$.
If $a=k\le\lambda m$, then
\[
 \frac{k-1}{m-k+1}
 \le\frac{2c+s}{cs-s-2c}
\]
and
\[
 \frac{t_0}{m-k+1}
 \le\frac{sc}{g(cs-s-2c)}.
\]
Using these estimates in~\eqref{eq:generic-order-bound} proves
\eqref{eq:k-error-bound}.
\end{proof}

For every index $h$, define
\begin{equation}\label{eq:S-h-definition}
 \Sset_h
 :=\{\Pp_j:|I_j\cap I_h|>\alpha\mathsf{AVG}\}.
\end{equation}
If $t_0<h\le m/s$ and $|\Sset_h|\ge m/c$, then
\eqref{eq:input-candidate-cost-identity} and
Lemma~\ref{lem:ordered-error-bounds}(i) imply
\begin{equation}\label{eq:R2-bound}
 \Obj(\Pp_h)
 \le R_2\OPT,
 \qquad
 R_2
 :=
 2+\frac{\beta(g-s)+s}{g(s-1)}-\frac{2\alpha}{c}.
\end{equation}

Suppose instead that $|\Sset_h|<m/c$.  Define
\begin{equation}\label{eq:N-h-definition}
 \N_h
 :=
 \left\{
 \Pp_j:
 h<j\le h+\frac{m}{c}+\frac{m}{s}
 \ \text{and}\ 
 |I_j\cap I_h|\le\alpha\mathsf{AVG}
 \right\}.
\end{equation}
The index interval in~\eqref{eq:N-h-definition} contains
$m/c+m/s$ input clusterings, and fewer than $m/c$ of them belong to
$\Sset_h$.  Consequently,
\begin{equation}\label{eq:N-h-size}
 |\N_h|\ge\frac{m}{s}.
\end{equation}

Fix $\Pp_k\in\N_h$ and define
\begin{equation}\label{eq:F-hk-definition}
 \F_{h,k}
 :=
 \left\{
 \Pp_j:
 |I_j\cap I_h|\le\alpha\mathsf{AVG}
 \ \text{and}\ 
 |I_j\cap I_k|\le\alpha\mathsf{AVG}
 \right\}.
\end{equation}
If $|\F_{h,k}|<m/s$, then every clustering outside $\F_{h,k}$ belongs to
$\Sset_h\cup\Sset_k$.  Therefore,
\begin{equation}\label{eq:S-k-size}
 |\Sset_k|
 >m\left(1-\frac1s-\frac1c\right).
\end{equation}
Moreover, because $h\le m/s$ and $k\in\N_h$,
\[
 k\le\frac{2c+s}{cs}m.
\]
Equations~\eqref{eq:input-candidate-cost-identity} and
\eqref{eq:k-error-bound} now give
\begin{equation}\label{eq:R3-bound}
 \Obj(\Pp_k)
 \le R_3\OPT,
\end{equation}
where
\begin{equation}\label{eq:R3-definition}
 R_3
 :=
 2+
 \frac{\beta\bigl(g(2c+s)-sc\bigr)+sc}
      {g(cs-s-2c)}
 -2\left(1-\frac1s-\frac1c\right)\alpha.
\end{equation}

Finally, if $|\F_{h,k}|\ge m/s$, then every
$\Pp_\ell\in\F_{h,k}$ satisfies
\[
 |I_h\cap I_k|\le\alpha\mathsf{AVG},\qquad
 |I_h\cap I_\ell|\le\alpha\mathsf{AVG},\qquad
 |I_k\cap I_\ell|\le\alpha\mathsf{AVG}.
\]
Hence Lemma~\ref{lem:sparse-triple} applies to the triple
$(h,k,\ell)$.

\subsubsection{Sampling the four cases}

For ease of exposition, assume that $m/g,m/s$, and $m/c$ are integers.
This is without loss of generality: one may replace every input clustering
by $12g$ identical copies.  This multiplies every candidate cost and
$\OPT$ by the same factor and does not affect any ratio.  Sampling a
uniformly random copy is equivalent to sampling a uniformly random original
input.  The replication need not be performed explicitly by the algorithm.

One \emph{macro-pool} consists of three independent multisets
\[
  \Sset^{(1)},\Sset^{(2)},\Sset^{(3)},
\]
each containing
\begin{equation}\label{eq:q0-definition}
  q_0:=\left\lceil8g\ln m\right\rceil
\end{equation}
indices sampled independently and uniformly from $[m]$.  The candidate set
of this macro-pool contains the input candidate for every sampled index and
the triple candidate for every ordered triple with one index from each of
$\Sset^{(1)},\Sset^{(2)},\Sset^{(3)}$.  Repetitions are harmless: they can be
viewed as distinct copies of the same input clustering.

\begin{lemma}[One-pool candidate lemma]\label{lem:one-pool}
At any fixed prefix, with probability at least $1-m^{-3}$, one macro-pool
contains a candidate of cost at most
\begin{equation}\label{eq:r-general}
 r\OPT,
 \qquad
 r:=\max\{R_1,R_2,R_3,R_4\},
\end{equation}
where
\begin{align}
 R_1&:=2-\beta,\label{eq:R1-definition}\\
 R_2&:=2+\frac{\beta(g-s)+s}{g(s-1)}-\frac{2\alpha}{c},
       \label{eq:R2-definition}\\
 R_3&:=2+
 \frac{\beta\bigl(g(2c+s)-sc\bigr)+sc}
      {g(cs-s-2c)}
 -2\left(1-\frac1s-\frac1c\right)\alpha,
       \label{eq:R3-definition-again}\\
 R_4&:=1+3(\eta+1)\alpha.\label{eq:R4-definition}
\end{align}
\end{lemma}

\begin{proof}
Let $t_0$ be defined by~\eqref{eq:low-inputs}.  First suppose
$t_0\ge m/g$.  Then at least an $1/g$ fraction of the input clusterings
satisfy $|I_i|\le(1-\beta)\mathsf{AVG}$.  The probability that
$\Sset^{(1)}$ misses all of them is at most
\[
 \left(1-\frac1g\right)^{q_0}
 \le e^{-q_0/g}
 \le m^{-8}.
\]
If $\Pp_i$ is sampled from this set, then
\eqref{eq:input-candidate-cost-identity} gives
\[
 \Obj(\Pp_i)
 \le\OPT+m(1-\beta)\mathsf{AVG}
 =(2-\beta)\OPT
 =R_1\OPT.
\]

It remains to consider $t_0<m/g$.  The set of indices
\[
 \{t_0+1,t_0+2,\ldots,m/s\}
\]
has size at least
\[
 \frac{m}{s}-\frac{m}{g}
 =\frac{g-s}{gs}m.
\]
Because $g\ge2s$, the probability that $\Sset^{(1)}$ misses this set is at
most
\[
 \exp\!\left(-q_0\frac{g-s}{gs}\right)
 \le m^{-8}.
\]
Condition on sampling one such index $h$.

If $|\Sset_h|\ge m/c$, then~\eqref{eq:R2-bound} shows that the sampled input
candidate $\Pp_h$ has cost at most $R_2\OPT$.  Otherwise,
\eqref{eq:N-h-size} shows that $|\N_h|\ge m/s$.  Conditioned on the already
chosen $h$, the independent multiset $\Sset^{(2)}$ misses $\N_h$ with
probability at most
\[
 \left(1-\frac1s\right)^{q_0}
 \le e^{-q_0/s}
 \le m^{-8}.
\]
Condition on sampling an index $k\in\N_h$.

If $|\F_{h,k}|<m/s$, then~\eqref{eq:R3-bound} shows that the sampled input
candidate $\Pp_k$ has cost at most $R_3\OPT$.  Otherwise,
$|\F_{h,k}|\ge m/s$, and the independent multiset $\Sset^{(3)}$ misses
$\F_{h,k}$ with probability at most $m^{-8}$.  If
$\ell\in\F_{h,k}$ is sampled, Lemma~\ref{lem:sparse-triple} shows that the
sampled triple candidate $(h,k,\ell)$ has cost at most $R_4\OPT$.

There are at most three relevant sampling failures in any branch of the
case analysis.  Their total probability is at most
\[
 3m^{-8}\le m^{-3}
\]
for $m\ge2$.  This proves the lemma.
\end{proof}

\subsubsection{Fixing the parameters}

We now fix
\begin{equation}\label{eq:parameter-choice}
 \alpha:=\frac{1}{6(\eta+1)},
 \qquad
 \beta:=\frac{\alpha}{24},
 \qquad
 s:=4,
 \qquad
 c:=3,
 \qquad
 g:=\left\lceil\frac{48}{\alpha}\right\rceil.
\end{equation}
Then $cs-s-2c=2>0$, $g>s$, and
\begin{equation}\label{eq:delta-definition}
 \delta:=\beta
 =\frac{1}{144(\eta+1)}.
\end{equation}

The four terms in Lemma~\ref{lem:one-pool} satisfy
\begin{align}
 R_1&=2-\frac{\alpha}{24}=2-\delta,
       \label{eq:R1-evaluated}\\
 R_2
 &=2+\frac{\beta(g-4)+4}{3g}-\frac{2\alpha}{3}
 \le2-\frac{5\alpha}{8},
       \label{eq:R2-evaluated}\\
 R_3
 &=2+5\beta+\frac{6-6\beta}{g}-\frac{5\alpha}{6}
 \le2-\frac{\alpha}{2},
       \label{eq:R3-evaluated}\\
 R_4&=1+3(\eta+1)\alpha=\frac32.
       \label{eq:R4-evaluated}
\end{align}
For completeness, the inequality in~\eqref{eq:R2-evaluated} uses
\[
 \frac{\beta}{3}=\frac{\alpha}{72}
 \qquad\text{and}\qquad
 \frac{4}{3g}\le\frac{\alpha}{36},
\]
while~\eqref{eq:R3-evaluated} uses
\[
 5\beta=\frac{5\alpha}{24}
 \qquad\text{and}\qquad
 \frac6g\le\frac{\alpha}{8}.
\]
Since $\eta\ge1$, $R_1$ is the largest of the four terms.  Therefore,
\begin{equation}\label{eq:r-subtwo}
  r=2-\delta.
\end{equation}

\subsection{The history-stability lemma}

In this subsection, we strengthen the pointwise sampled-candidate lemma to a
statement about the complete history of a candidate.  This is the central
new ingredient needed for rebuilding.

Let $\A$ denote the union of all candidate processes maintained by the
algorithm, and define
\begin{equation}\label{eq:M-definition}
  M_t:=\min_{a\in\A}G_a(t).
\end{equation}
Because every $G_a(t)$ is nondecreasing, $M_t$ is also nondecreasing.

\begin{lemma}[History stability]\label{lem:history-stability}
Fix a prefix $t$.  On the sampling-success event of
Lemma~\ref{lem:one-pool}, the corresponding macro-pool contains a candidate
$a$ satisfying
\begin{equation}\label{eq:history-good-candidate}
  G_a(t)\le(2-\delta)\OPT_t.
\end{equation}
\end{lemma}

\begin{proof}
Fix an optimal clustering $\F_t^*$ of $V_t$.  Define all error sets at this
prefix relative to this single clustering $\F_t^*$.  In particular,
\begin{equation}\label{eq:time-t-error-set}
 I_i^t
 :=
 \left\{
  \{x,y\}\subseteq V_t:
  \Pp_i^t\text{ and }\F_t^*
  \text{ disagree on }\{x,y\}
 \right\}.
\end{equation}

In the first three cases of Lemma~\ref{lem:one-pool}, the good candidate is
an input candidate $i$.  Its score is nondecreasing, and therefore
\[
 G_i(t)=f_i(t)\le(2-\delta)\OPT_t.
\]

It remains to consider the sparse-triple case.  Suppose the sampled triple
$\tau=(i,j,k)$ satisfies
\[
 |I_i^t\cap I_j^t|,
 |I_j^t\cap I_k^t|,
 |I_k^t\cap I_i^t|
 \le\alpha\frac{\OPT_t}{m}.
\]
Define
\begin{equation}\label{eq:Btt-definition}
 B_{t,t}
 :=
 (I_i^t\cap I_j^t)
 \cup(I_j^t\cap I_k^t)
 \cup(I_k^t\cap I_i^t).
\end{equation}
For an earlier time $u\le t$, let
\begin{equation}\label{eq:Btu-definition}
 B_{t,u}:=B_{t,t}\cap\binom{V_u}{2}.
\end{equation}
It is important that both $B_{t,t}$ and $B_{t,u}$ are defined relative to
the same final-prefix optimum $\F_t^*$.  In general, the analogous error
sets defined using unrelated optima $\F_u^*$ and $\F_t^*$ need not be
nested.

The majority graph $H_\tau(u)$ is the restriction of $H_\tau(t)$ to $V_u$.
Consequently, $B_{t,u}$ is exactly the set of edges on which the majority
sign in $H_\tau(u)$ disagrees with the pair relation of
$\F_t^*|_{V_u}$.  Therefore,
\begin{equation}\label{eq:restricted-opt-majority-error}
 \err_{H_\tau(u)}(\F_t^*|_{V_u})=|B_{t,u}|.
\end{equation}
The Online Agreement guarantee gives
\begin{align}
 \err_{H_\tau(u)}(\A_\tau(u))
 &\le\eta\min_{\C}\err_{H_\tau(u)}(\C)\notag\\
 &\le\eta|B_{t,u}|.                       \label{eq:agreement-history-bound}
\end{align}
Using the Hamming triangle inequality through the majority sign vector,
we get
\begin{equation}\label{eq:history-distance-bound}
 \dist_{V_u}(\A_\tau(u),\F_t^*|_{V_u})
 \le(\eta+1)|B_{t,u}|.
\end{equation}

We now compare the full consensus cost at time $u$ to the optimum at the
later time $t$.  By the triangle inequality and restriction monotonicity,
\begin{align}
 f_\tau(u)
 &=\sum_{h=1}^m\dist_{V_u}(\A_\tau(u),\Pp_h^u)\notag\\
 &\le
 \sum_{h=1}^m\dist_{V_u}(\Pp_h^u,\F_t^*|_{V_u})
 +m\dist_{V_u}(\F_t^*|_{V_u},\A_\tau(u))\notag\\
 &\le
 \OPT_t+m(\eta+1)|B_{t,u}|\notag\\
 &\le
 \OPT_t+m(\eta+1)|B_{t,t}|\notag\\
 &\le
 \bigl(1+3(\eta+1)\alpha\bigr)\OPT_t\notag\\
 &=\frac32\OPT_t.                         \label{eq:history-triple-final}
\end{align}
This bound holds for every $u\le t$.  Taking the maximum over $u$ gives
\[
 G_\tau(t)\le\frac32\OPT_t\le(2-\delta)\OPT_t.
\]
This completes the proof.
\end{proof}

\subsection{Amplification over all prefixes}

In this subsection, we make the history-stability guarantee simultaneous
over all $n$ prefixes.  Let
\begin{equation}\label{eq:Q-definition}
 Q
 :=
 \max\left\{
  1,
  \left\lceil\frac{\ln(n/\zeta)}{3\ln m}\right\rceil
 \right\}.
\end{equation}
At initialization, independently construct $Q$ macro-pools.  The final
candidate family $\A$ is the union of the candidate families generated by
these macro-pools.  Let
\begin{equation}\label{eq:K-definition}
 K:=|\A|.
\end{equation}
Since each macro-pool contains at most $3q_0$ input candidates and at most
$q_0^3$ triple candidates,
\begin{equation}\label{eq:K-bound}
 K\le Q(3q_0+q_0^3)
 =O_\eta\!\left(\log^3\!\frac{mn}{\zeta}\right).
\end{equation}

For a fixed prefix, one macro-pool fails with probability at most $m^{-3}$.
Because the $Q$ macro-pools are independent, they all fail at this prefix
with probability at most
\[
 m^{-3Q}\le\frac{\zeta}{n}.
\]
A union bound over $t=1,2,\ldots,n$ proves the following event.

\begin{lemma}[Simultaneous history event]\label{lem:simultaneous-event}
With probability at least $1-\zeta$, simultaneously for every time
$t\in[n]$,
\begin{equation}\label{eq:simultaneous-M-bound}
 M_t\le(2-\delta)\OPT_t.
\end{equation}
\end{lemma}

No independence between different prefixes is required for the final union
bound.  Independence is used only between the $Q$ macro-pools at a fixed
prefix.

\subsection{The rebuilding algorithm}

In this subsection, we describe the rebuilding rule.  Define the hysteresis
parameter
\begin{equation}\label{eq:theta-definition}
  \theta:=\frac{\delta}{4}.
\end{equation}
The algorithm maintains one current candidate.  Its output clustering is
the current clustering of this candidate.

\begin{figure}[H]
\centering
\fbox{\begin{minipage}{0.94\textwidth}
\textbf{Algorithm 1: Sampled cluster fitting with rebuilding}
\begin{enumerate}[leftmargin=2em,itemsep=3pt,topsep=5pt]
\item Before the first arrival, construct the $Q$ independent macro-pools
      and initialize every input and triple candidate described above.
\item At time $t$, update every candidate $\A_a(t)$, its current cost
      $f_a(t)$, its history score $G_a(t)$, and
      $M_t=\min_aG_a(t)$.
\item If $M_t=0$, output any candidate with history score zero.  If the
      current candidate becomes positive while $M_t$ remains zero, switch to
      another zero-score candidate.
\item When $M_t$ first becomes positive, or when the preceding positive
      phase ends, begin a new phase.  Set the phase baseline to $B:=M_t$,
      clear the phase's discarded-candidate set, and choose a minimizer of
      $G_a(t)$ as the current candidate.
\item A positive phase continues while $M_t<(1+\theta)B$.  Within the phase,
      follow the current candidate while
      $G_a(t)<(1+\theta)B$.  If this inequality fails, discard $a$ for the
      remainder of the phase and switch to any undiscarded candidate whose
      history score is below $(1+\theta)B$.
\item At the first time $M_t\ge(1+\theta)B$, end the current phase and start
      a new phase using Step~4.
\end{enumerate}
\end{minipage}}
\caption{The rebuilding algorithm.}\label{alg:rebuilding}
\end{figure}

The switch in Step~5 is always possible.  Indeed, while the current phase
continues, the candidate attaining $M_t$ has score strictly smaller than
$(1+\theta)B$.  A candidate discarded earlier in the phase has score at
least $(1+\theta)B$ forever because history scores are nondecreasing.  Thus,
the current minimizer has not been discarded.

\subsection{Analysis of the rebuilding algorithm}

In this subsection, we prove the competitive ratio and the per-vertex
recourse bound.

\subsubsection{Competitive ratio}

\begin{lemma}[Cost within a phase]\label{lem:phase-cost}
Assume the simultaneous history event in
Lemma~\ref{lem:simultaneous-event}.  At every time $t$ in a positive phase,
the output of Algorithm~\ref{alg:rebuilding} satisfies
\begin{equation}\label{eq:phase-cost-bound}
 \Obj_t(\C_t)
 <(1+\theta)(2-\delta)\OPT_t.
\end{equation}
If $M_t=0$, then the output cost is zero.
\end{lemma}

\begin{proof}
If $M_t=0$, the algorithm follows a candidate with history score zero.
Its current cost is at most its history score, so the output cost is zero.

Consider a positive phase beginning at time $t_0$ with baseline
$B=M_{t_0}$.  Throughout the phase, the maintained candidate $a$ satisfies
\[
 f_a(t)\le G_a(t)<(1+\theta)B.
\]
By the simultaneous event,
\[
 B=M_{t_0}\le(2-\delta)\OPT_{t_0}.
\]
The monotonicity of the optimum gives $\OPT_{t_0}\le\OPT_t$.  Consequently,
\[
 \Obj_t(\C_t)=f_a(t)
 <(1+\theta)(2-\delta)\OPT_t.
\]
At a phase boundary, the algorithm first starts the new phase and chooses a
minimizer of $G_a(t)$, so the same bound holds there as well.
\end{proof}

Using $\theta=\delta/4$, we have
\begin{align}
 (1+\theta)(2-\delta)
 &=\left(1+\frac\delta4\right)(2-\delta)\notag\\
 &=2-\frac\delta2-\frac{\delta^2}{4}\notag\\
 &\le2-\frac{1}{288(\eta+1)}.             \label{eq:final-factor-arithmetic}
\end{align}
Combining Lemma~\ref{lem:phase-cost},
Lemma~\ref{lem:simultaneous-event}, and
\eqref{eq:final-factor-arithmetic} proves the competitive-ratio part of
Theorem~\ref{thm:main}.

\subsubsection{Number of rebuilds}

Every candidate cost is an integer between zero and
\begin{equation}\label{eq:Lambda-definition}
 \Lambda:=m\binom n2.
\end{equation}
The same is true of every history score and of $M_t$.  During the zero-score
stage, a candidate that becomes positive can never again have score zero.
Hence there are at most $K$ candidate switches while $M_t=0$.

Let $B_1,B_2,\ldots,B_H$ be the baselines of the positive phases.  By the
phase-ending rule,
\[
 B_{h+1}\ge(1+\theta)B_h.
\]
Since $1\le B_h\le\Lambda$, the number of positive phases is at most
\begin{equation}\label{eq:number-phases}
 H
 \le
 1+\left\lceil\log_{1+\theta}\Lambda\right\rceil
 =O\!\left(\frac{1}{\delta}\log(mn)\right).
\end{equation}
Within one phase, every discarded candidate remains discarded until the
phase ends.  Therefore there are at most $K$ candidate switches in each
phase, up to one additional choice at the phase boundary.  The total number
of outer candidate switches is consequently
\begin{equation}\label{eq:outer-switches}
 O\!\left(
 K\left[1+\frac{1}{\delta}\log(mn)\right]
 \right).
\end{equation}

\subsubsection{Per-vertex recourse}

Fix a vertex $v$.  Whenever the outer algorithm switches from one candidate
to another, the cluster identifier assigned to $v$ changes at most once.
Thus~\eqref{eq:outer-switches} bounds the recourse caused by outer rebuilding.

While an input candidate is selected, it causes no internal recourse.  For a
fixed triple candidate $\tau$, Theorem~\ref{thm:online-agreement} bounds the
number of changes of $v$ inside that candidate by
$q_{\mathrm A}\lceil\log_2 n\rceil$ over the entire arrival sequence.  If
the outer algorithm leaves $\tau$ and later returns to it, the re-entry is
already counted as an outer candidate switch.  Therefore, summing
conservatively over all $K$ candidates,
\begin{equation}\label{eq:recourse-explicit}
 \rec(v)
 \le
 O\!\left(
 K\left[
   1+\frac{1}{\delta}\log(mn)
   +q_{\mathrm A}\log n
 \right]
 \right).
\end{equation}
Using~\eqref{eq:K-bound} and the fact that $\eta,q_{\mathrm A}$ are absolute
constants gives
\begin{equation}\label{eq:recourse-polylog}
 \rec(v)
 =O_{\eta,q_{\mathrm A}}\!\left(
    \log^4\!\frac{mn}{\zeta}
  \right).
\end{equation}
This bound holds for every outcome of the random samples.  Together with the
competitive analysis, it completes the proof of
Theorem~\ref{thm:main}.

\subsection{Discussion and limitations}

In this subsection, we record the scope of the result and explain why the
cluster-fitting candidates are necessary.

\begin{remark}[Why the previous sampled-medoid argument stops at $2$]
For an input candidate $\Pp_i$, define
\[
 F_i(t):=\sum_{j=1}^m\dist_{V_t}(\Pp_i^t,\Pp_j^t).
\]
A uniformly random input candidate satisfies
\[
 \mathbb E_i[F_i(t)]
 \le2\left(1-\frac1m\right)\OPT_t.
\]
For a target ratio $R$, Markov's inequality guarantees an $R$-good fraction
only when
\[
 R>2\left(1-\frac1m\right).
\]
Thus, for a fixed constant $R<2$ and arbitrary $m$, sampling input
clusterings alone gives no useful good-candidate probability.  The sparse
triple candidate avoids this obstruction by providing the stronger
$3\OPT/2$ branch.
\end{remark}

\begin{remark}[Oblivious adversary]
The probability amplification treats every prefix as a fixed instance and
then takes a union bound.  Therefore, the input clusterings and arrival order
must be chosen independently of the initial random samples.  The proof does
not cover an adversary that observes the algorithm's outputs, infers
information about the samples, and adaptively chooses future vertices.
\end{remark}

\begin{remark}[Unknown horizon]
The proof assumes that an upper bound on $n$ is known when the sample pools
are chosen.  If the horizon is unknown, one may use geometrically increasing
horizons and restart the outer rebuilding mechanism when the current horizon
is exceeded.  Each vertex participates in only logarithmically many such
restarts, so this changes the recourse by at most another logarithmic factor.
\end{remark}

\begin{remark}[Deterministic version]
Enumerating all $m$ input candidates and all $\binom m3$ triple candidates
removes the sampling failure and gives the same competitive analysis
deterministically.  However, the rebuilding proof then gives polynomial,
rather than polylogarithmic, dependence on $m$ in the recourse.  Obtaining a
polylogarithmic-size deterministic hitting family for the four-case
cluster-fitting argument remains the missing ingredient for a deterministic
single-digit, or deterministic sub-$2$, result by this route.
\end{remark}

\begin{remark}[Computation]
The algorithm evaluates every maintained candidate against all $m$ input
clusterings.  This is polynomial time, but the present proof optimizes the
competitive ratio and recourse rather than update time or memory.  No random
evaluation sample is needed, and therefore the evaluation loss appearing in
the streaming-over-input-clusterings algorithm of
\cite{ChakrabortyEtAl2026} is absent here.
\end{remark}

\begin{thebibliography}{9}

\bibitem{ChakrabortyEtAl2026}
Diptarka Chakraborty, Kushagra Chatterjee, Debarati Das, and
Tien-Long Nguyen.
\newblock A generic framework for fair consensus clustering in streams.
\newblock In \emph{Proceedings of the 25th International Conference on
Autonomous Agents and Multiagent Systems (AAMAS)}, 2026.

\bibitem{CohenAddadEtAl2022}
Vincent Cohen-Addad, Silvio Lattanzi, Andreas Maggiori, and Nikos Parotsidis.
\newblock Online and consistent correlation clustering.
\newblock In \emph{Proceedings of the 39th International Conference on
Machine Learning (ICML)}, 2022.

\end{thebibliography}

\end{document}
